\documentclass[11pt]{article}

% ============================================================
% Paquetes
% ============================================================
\usepackage[utf8]{inputenc}
\usepackage[T1]{fontenc}
\usepackage[spanish]{babel}
\usepackage[margin=2.5cm]{geometry}
\usepackage{amsmath,amssymb,amsthm}
\usepackage{tikz}
\usetikzlibrary{babel}
\usepackage{xcolor}
\usepackage{listings}
\usepackage{hyperref}

% ------------------------------------------------------------
% Entorno para comentarios explicativos / observaciones
% ------------------------------------------------------------
\theoremstyle{remark}
\newtheorem*{nota}{Nota}

% ------------------------------------------------------------
% Estilo para fragmentos de código Python
% ------------------------------------------------------------
\lstset{
  language=Python,
  basicstyle=\ttfamily\small,
  keywordstyle=\color{blue!70!black}\bfseries,
  commentstyle=\color{gray}\itshape,
  stringstyle=\color{red!60!black},
  showstringspaces=false,
  breaklines=true,
  frame=single,
  framerule=0.4pt,
  numbers=left,
  numbersep=8pt,
  numberstyle=\tiny\color{gray},
  xleftmargin=15pt
}

\hypersetup{
  colorlinks=true,
  linkcolor=blue!50!black,
  urlcolor=blue!50!black
}

% ============================================================
\title{\textbf{Inferencia Bayesiana en Regresión Lineal Simple}\\[4pt]
\large Verosimilitud, distribución previa y estimación MAP\\
\normalsize (Continuación de la clase anterior)}
\author{Apuntes de clase --- Estadística}
\date{}
% ============================================================

\begin{document}

\maketitle

\begin{abstract}
\noindent
Estas notas retoman el análisis bayesiano del modelo de regresión lineal
simple $y_i = a + bx_i + \varepsilon_i$. Se construye la función de
verosimilitud, se especifica una distribución previa no informativa para los
parámetros $(a,b)$, y se combinan ambas mediante el teorema de Bayes para
obtener la densidad posterior. Finalmente se deduce el estimador de Máximo a
Posteriori (MAP) resolviendo un sistema lineal $M\vec\theta=\vec Z$, y se
muestra que la posterior completa de $(a,b)$ es, de hecho, una distribución
normal bivariada.
\end{abstract}

\tableofcontents

% =====================================================
\section{Modelo y función de verosimilitud}
% =====================================================

\subsection{Especificación del modelo}

Se considera el modelo de regresión lineal simple
\begin{equation}
y_i = a + bx_i + \varepsilon_i, \qquad i=1,\dots,n,
\label{eq:modelo}
\end{equation}
donde $a$ (intercepto) y $b$ (pendiente) son los \textbf{parámetros} del
modelo, y los errores se suponen independientes e idénticamente distribuidos:
\begin{equation}
\varepsilon_i \sim \mathcal{N}(0,\sigma^2), \qquad \sigma = 0.3 \ \ (\text{conocido}).
\label{eq:error}
\end{equation}

\begin{nota}
A diferencia de $a$ y $b$, que son desconocidos y sobre los cuales se hará
toda la inferencia, $\sigma$ se trata aquí como una \emph{constante conocida},
fijada en $0.3$. Por ello, en lo que sigue $\theta=(a,b)$ denota el vector de
parámetros desconocidos del modelo.
\end{nota}

\subsection{Distribución de cada observación}

Escribimos $M(a,b)$ para referirnos al modelo evaluado en $x_i$, es decir
$M(a,b)\equiv a+bx_i$. Como $y_i$ es una traslación de $\varepsilon_i$
(ecuación \eqref{eq:modelo}) y $\varepsilon_i\sim\mathcal N(0,\sigma^2)$, se
sigue que
\begin{equation}
y_i \mid M(a,b) \ \sim\ \mathcal{N}\!\big(a+bx_i,\ \sigma^2\big).
\label{eq:dist-yi}
\end{equation}

Por lo tanto, el \emph{núcleo} de la densidad de cada observación (es decir,
la densidad sin su constante de normalización, que no depende de $a,b$) es
\begin{equation}
f(y_i\mid M(a,b)) \ \propto\ \exp\!\left(-\frac{\big(y_i-(a+bx_i)\big)^2}{2\sigma^2}\right).
\label{eq:fyi}
\end{equation}

\subsection{Verosimilitud conjunta}

Si las observaciones $y_1,\dots,y_n$ son condicionalmente independientes dado
el modelo, la verosimilitud conjunta es el producto de las densidades
individuales:
\begin{equation}
L(y\mid M(a,b)) = \prod_{i=1}^{n} f(y_i\mid M(a,b)).
\label{eq:L-producto}
\end{equation}

Como el producto de exponenciales es la exponencial de la suma de los
exponentes, sustituyendo \eqref{eq:fyi} en \eqref{eq:L-producto}:
\begin{equation}
L(y\mid M(a,b)) \ \propto\ \exp\!\left(-\frac{1}{2\sigma^2}\sum_{i=1}^{n}\big(y_i-(a+bx_i)\big)^2\right).
\label{eq:L-final}
\end{equation}

% =====================================================
\section{Distribución previa}
% =====================================================

\subsection{Prior no informativo para \texorpdfstring{$a$ y $b$}{a y b}}

A cada parámetro se le asigna una distribución previa normal centrada en
cero, con una desviación estándar grande, es decir, poco informativa, ya que
asigna probabilidad apreciable a un rango muy amplio de valores:
\begin{equation}
a \sim \mathcal{N}(0,\sigma_a^2), \qquad b\sim\mathcal N(0,\sigma_b^2),
\qquad \sigma_a=\sigma_b=10.
\label{eq:prior-ab}
\end{equation}

Los núcleos de estas densidades son
\begin{equation}
g(a)\propto\exp\!\left(-\frac{a^2}{2\sigma_a^2}\right), \qquad
g(b)\propto\exp\!\left(-\frac{b^2}{2\sigma_b^2}\right).
\label{eq:g-ab}
\end{equation}

\subsection{Distribución previa conjunta}

Suponiendo independencia a priori entre $a$ y $b$, la densidad previa
conjunta es el producto de las marginales:
\begin{equation}
f(a,b) = g(a)\,g(b) \ \propto\ \exp\!\left(-\frac{a^2}{2\sigma_a^2}\right)
\exp\!\left(-\frac{b^2}{2\sigma_b^2}\right).
\label{eq:f-ab}
\end{equation}

Como $\sigma_a=\sigma_b=10$, se tiene $2\sigma_a^2=2\sigma_b^2=2\times10^2=200$,
de modo que \eqref{eq:f-ab} se simplifica a
\begin{equation}
f(a,b)\ \propto\ \exp\!\left(-\frac{a^2+b^2}{200}\right).
\label{eq:f-ab-final}
\end{equation}

% =====================================================
\section{Distribución posterior}
% =====================================================

\subsection{Teorema de Bayes}

La densidad posterior conjunta de $(a,b)$ dados los datos $y=(y_1,\dots,y_n)$
es proporcional al producto de la verosimilitud \eqref{eq:L-final} y la
previa \eqref{eq:f-ab-final}:
\begin{equation}
P(M(a,b)\mid y) \ \propto\ L(y\mid M(a,b))\cdot f(a,b).
\label{eq:bayes}
\end{equation}

\subsection{Núcleo de la posterior}

Sustituyendo \eqref{eq:L-final} y \eqref{eq:f-ab-final} en \eqref{eq:bayes}:
\begin{equation}
P(a,b\mid y)\ \propto\
\exp\!\left(-\frac{1}{2\sigma^2}\sum_{i}(y_i-a-bx_i)^2\right)
\cdot
\exp\!\left(-\frac{a^2+b^2}{200}\right),
\label{eq:posterior1}
\end{equation}
o, de forma equivalente, agrupando ambos factores en un solo exponente:
\begin{equation}
P(a,b\mid y)\ \propto\
\exp\!\left(-\frac{1}{2\sigma^2}\sum_{i}(y_i-a-bx_i)^2 -\frac{a^2+b^2}{200}\right).
\label{eq:posterior2}
\end{equation}

\subsection{Probabilidades posteriores}

Una vez conocida (al menos como núcleo) la densidad posterior de un parámetro
$\theta$, las probabilidades de interés se obtienen como áreas bajo dicha
densidad. Por ejemplo, la probabilidad de que $\theta$ supere un valor de
referencia $\theta_1$ es
\begin{equation}
P(\theta>\theta_1\mid x) = \int_{\theta_1}^{\infty} f(\theta\mid x)\,d\theta,
\label{eq:prob-posterior}
\end{equation}
que corresponde gráficamente al área sombreada de la Figura~\ref{fig:posterior}.

\begin{figure}[h]
\centering
\begin{tikzpicture}
  % Curva de densidad completa
  \draw[thick, blue!60!black, domain=-3.3:3.3, samples=120, smooth]
    plot ({\x},{2.3*exp(-\x*\x/2)});
  % Region sombreada bajo la curva, desde theta_1 hacia la derecha
  \fill[blue!15] plot[domain=1.2:3.3, samples=60, smooth]
    ({\x},{2.3*exp(-\x*\x/2)}) -- (3.3,0) -- (1.2,0) -- cycle;
  % Redibuja la curva sobre la region sombreada
  \draw[thick, blue!60!black, domain=1.2:3.3, samples=60, smooth]
    plot ({\x},{2.3*exp(-\x*\x/2)});
  % Eje horizontal
  \draw[->] (-3.4,0) -- (3.6,0) node[right] {$\theta$};
  % Lineas de referencia verticales
  \draw[dashed] (0,0) -- (0,2.35);
  \draw[dashed] (1.2,0) -- (1.2,{2.3*exp(-1.2*1.2/2)});
  % Etiquetas sobre el eje
  \node[below] at (0,-0.05) {$\bar\theta$};
  \node[below] at (1.2,-0.05) {$\theta_1$};
  % Etiqueta de la probabilidad
  \node at (2.35,0.55) {$P(\theta>\theta_1\mid x)$};
\end{tikzpicture}
\caption{Densidad posterior de un parámetro $\theta$. El área sombreada a la
derecha de $\theta_1$ representa $P(\theta>\theta_1\mid x)$, mientras que
$\bar\theta$ indica la media (o moda) posterior.}
\label{fig:posterior}
\end{figure}

% =====================================================
\section{Estimación MAP}
\label{sec:map}
% =====================================================

\subsection{Logaritmo de la posterior}

Para encontrar el estimador de Máximo a Posteriori (MAP) conviene trabajar
con el logaritmo de la posterior: el logaritmo es una función estrictamente
creciente, por lo que no cambia la ubicación del máximo, y transforma el
producto \eqref{eq:posterior1} en una suma, mucho más fácil de derivar.
\begin{equation}
\ell(\theta\mid y) = \ln(\text{Post}).
\label{eq:loglik-def}
\end{equation}

Tomando el logaritmo del núcleo \eqref{eq:posterior2}:
\begin{equation}
\ell(\theta\mid y) = -\frac{1}{2\sigma^2}\sum_{i=1}^{n}(y_i-a-bx_i)^2
- \frac{a^2+b^2}{200}.
\label{eq:loglik}
\end{equation}

\subsection{Condiciones de primer orden}

El estimador MAP $(\hat a,\hat b)$ es el punto donde $\ell(\theta\mid y)$
alcanza su máximo, es decir, donde sus derivadas parciales se anulan.

\paragraph{Derivada respecto de $a$.}
\begin{equation}
\frac{\partial \ell}{\partial a}
= \frac{1}{\sigma^2}\sum_{i=1}^n (y_i-a-bx_i) - \frac{a}{100}.
\label{eq:dl-da}
\end{equation}
Igualando a cero en $(\hat a,\hat b)$:
\begin{equation}
\frac{1}{\sigma^2}\sum_{i=1}^n (y_i-\hat a-\hat b x_i) - \frac{\hat a}{100} = 0.
\label{eq:foc-a}
\end{equation}

\paragraph{Derivada respecto de $b$.}
\begin{equation}
\frac{\partial \ell}{\partial b}
= \frac{1}{\sigma^2}\sum_{i=1}^n x_i(y_i-a-bx_i) - \frac{b}{100}.
\label{eq:dl-db}
\end{equation}
Igualando a cero:
\begin{equation}
\frac{1}{\sigma^2}\sum_{i=1}^n x_i(y_i-\hat a-\hat b x_i) - \frac{\hat b}{100} = 0.
\label{eq:foc-b}
\end{equation}

\begin{nota}
El factor $1/100$ que aparece en \eqref{eq:foc-a} y \eqref{eq:foc-b} proviene
de derivar el término de la previa, $-\dfrac{a^2+b^2}{200}$: como
$\dfrac{\partial}{\partial a}\!\left(-\dfrac{a^2}{200}\right)=-\dfrac{2a}{200}
=-\dfrac{a}{100}$ (y análogamente para $b$), el $2$ del numerador se cancela
con el $2$ del denominador $200$, dejando $1/100$.
\end{nota}

\subsection{Sistema lineal y forma matricial}

Las ecuaciones \eqref{eq:foc-a} y \eqref{eq:foc-b} son \textbf{lineales} en
$\hat a$ y $\hat b$. Multiplicando ambas por $\sigma^2$, distribuyendo las
sumatorias y agrupando los términos en $\hat a$ y $\hat b$ a la izquierda:

De \eqref{eq:foc-a}:
\begin{equation}
\left(n+\frac{\sigma^2}{100}\right)\hat a + \left(\sum_{i}x_i\right)\hat b
= \sum_i y_i.
\label{eq:eq1}
\end{equation}

De \eqref{eq:foc-b}:
\begin{equation}
\left(\sum_i x_i\right)\hat a
+ \left(\sum_i x_i^2+\frac{\sigma^2}{100}\right)\hat b
= \sum_i x_iy_i.
\label{eq:eq2}
\end{equation}

Este sistema de dos ecuaciones lineales en $(\hat a,\hat b)$ se escribe en
forma matricial $M\vec\theta=\vec Z$ como
\begin{equation}
\underbrace{\begin{bmatrix}
n+\dfrac{\sigma^2}{100} & \displaystyle\sum_i x_i \\[10pt]
\displaystyle\sum_i x_i & \displaystyle\sum_i x_i^2+\dfrac{\sigma^2}{100}
\end{bmatrix}}_{M}
\underbrace{\begin{bmatrix}\hat a \\ \hat b\end{bmatrix}}_{\vec\theta}
=
\underbrace{\begin{bmatrix}
\displaystyle\sum_i y_i \\[6pt] \displaystyle\sum_i x_iy_i
\end{bmatrix}}_{\vec Z}.
\label{eq:sistema-matricial}
\end{equation}

\begin{nota}[Interpretación: regularización tipo \emph{ridge}]
Si en \eqref{eq:sistema-matricial} se ignoraran los términos $\sigma^2/100$
(lo que equivaldría a $\sigma_a,\sigma_b\to\infty$, es decir, una previa
totalmente plana), la matriz $M$ se reduciría a la matriz normal $X^\top X$
de mínimos cuadrados ordinarios, y $\vec Z$ sería $X^\top y$. Los términos
$\sigma^2/100$ añaden una penalización a valores grandes de $a$ y $b$
—exactamente el mecanismo de la regresión \emph{ridge}— con parámetro de
regularización efectivo $\lambda=\sigma^2/100$.
\end{nota}

% =====================================================
\section{Solución computacional del sistema}
% =====================================================

El sistema \eqref{eq:sistema-matricial} puede resolverse en Python de manera
\textbf{numérica} (con \texttt{NumPy}) o \textbf{simbólica} (con
\texttt{SymPy}).

\subsection{Solución numérica con NumPy}

\begin{lstlisting}
import numpy as np

sigma = 0.3

# n, sum_x, sum_x2, sum_y, sum_xy: estadisticos calculados a partir de los datos
M = np.array([
    [n + sigma**2 / 100,        sum_x],
    [sum_x,                     sum_x2 + sigma**2 / 100]
])
Z = np.array([sum_y, sum_xy])

a_map, b_map = np.linalg.solve(M, Z)
\end{lstlisting}


\subsection{Solución simbólica con SymPy}

\begin{lstlisting}
import sympy as sp

a, b = sp.symbols('a b')
n, sigma = sp.symbols('n sigma', positive=True)
sum_x, sum_x2, sum_y, sum_xy = sp.symbols('sum_x sum_x2 sum_y sum_xy')

eq1 = sp.Eq((n + sigma**2/100) * a + sum_x * b, sum_y)
eq2 = sp.Eq(sum_x * a + (sum_x2 + sigma**2/100) * b, sum_xy)

solucion = sp.solve([eq1, eq2], [a, b])
a_map_expr = solucion[a]
b_map_expr = solucion[b]
\end{lstlisting}

\begin{nota}
La ventaja de resolver el sistema simbólicamente es obtener $\hat a$ y
$\hat b$ como expresiones en función de $n,\sigma,\sum x_i,\sum x_i^2,\sum y_i$
y $\sum x_iy_i$. Esto permite, por ejemplo, verificar que en el límite
$\sigma_a,\sigma_b\to\infty$ (previa "plana") se recupera el estimador de
mínimos cuadrados ordinarios.
\end{nota}

% =====================================================
\section{La posterior completa: conexión con la normal bivariada}
% =====================================================

Hasta ahora se obtuvo únicamente el punto $(\hat a,\hat b)$ que maximiza la
posterior (su \emph{moda}, o MAP). Sin embargo, la expresión
\eqref{eq:posterior2} permite caracterizar \textbf{toda} la distribución
posterior de $(a,b)$, no solo su moda.

\subsection{Forma general de la normal bivariada}

La densidad normal bivariada tiene un núcleo de la forma (ver, por ejemplo,
Devore, capítulo de distribuciones conjuntas continuas):
\begin{equation}
f(x,y)\ \propto\ \exp\!\left(-\frac12\Big(Ax^2+2Bxy+Cy^2-2Dx-2Ey+F\Big)\right).
\label{eq:bivariada-general}
\end{equation}

\subsection{Identificación de los coeficientes}

Expandiendo el cuadrado en \eqref{eq:posterior2}:
\begin{equation}
\sum_i(y_i-a-bx_i)^2 = \sum_i y_i^2 - 2a\sum_i y_i - 2b\sum_i x_iy_i
+ n\,a^2 + 2ab\sum_i x_i + b^2\sum_i x_i^2,
\label{eq:expansion}
\end{equation}
y sustituyendo en \eqref{eq:posterior2}, agrupando los términos en $a^2$,
$ab$, $b^2$, $a$, $b$ y las constantes, el exponente de la posterior toma la
forma \eqref{eq:bivariada-general} con $(x,y)\to(a,b)$ y
\begin{align}
A &= \frac{n}{\sigma^2}+\frac{1}{100}, \label{eq:coefA}\\
B &= \frac{\sum_i x_i}{\sigma^2}, \label{eq:coefB}\\
C &= \frac{\sum_i x_i^2}{\sigma^2}+\frac{1}{100}, \label{eq:coefC}\\
D &= \frac{\sum_i y_i}{\sigma^2}, \label{eq:coefD}\\
E &= \frac{\sum_i x_iy_i}{\sigma^2}. \label{eq:coefE}
\end{align}

\begin{nota}
Comparando \eqref{eq:coefA}--\eqref{eq:coefC} con la matriz $M$ de
\eqref{eq:sistema-matricial}, se observa que
\[
M=\sigma^2\begin{bmatrix}A&B\\B&C\end{bmatrix},
\]
es decir, el sistema $M\vec\theta=\vec Z$ es simplemente el sistema
$\begin{bmatrix}A&B\\B&C\end{bmatrix}\vec\theta=\begin{bmatrix}D\\E\end{bmatrix}$
multiplicado por $\sigma^2$ en ambos lados (lo cual no altera la solución).
\end{nota}

\subsection{Matriz de precisión y matriz de covarianza posterior}

En la forma \eqref{eq:bivariada-general}, la matriz de coeficientes
cuadráticos es la \textbf{matriz de precisión} (la inversa de la matriz de
covarianza) de la distribución:
\begin{equation}
\Sigma^{-1} = \begin{bmatrix} A & B \\ B & C\end{bmatrix}.
\label{eq:precision}
\end{equation}

Invirtiendo esta matriz $2\times2$ se obtiene la \textbf{matriz de
covarianza} de la posterior:
\begin{equation}
\Sigma = \frac{1}{AC-B^2}\begin{bmatrix} C & -B \\ -B & A\end{bmatrix},
\label{eq:covarianza}
\end{equation}
con $A$, $B$ y $C$ dados por \eqref{eq:coefA}--\eqref{eq:coefC}.

\begin{nota}[Resultado principal]
Como tanto la verosimilitud \eqref{eq:L-final} como la previa
\eqref{eq:f-ab-final} son gaussianas, la posterior también lo es: este es un
caso de \textbf{previa conjugada normal-normal}.
\begin{equation}
(a,b)\mid y \ \sim\ \mathcal{N}_2\big((\hat a,\hat b),\ \Sigma\big),
\label{eq:posterior-bivariada}
\end{equation}
donde $(\hat a,\hat b)$ es la solución de \eqref{eq:sistema-matricial} (para
una distribución normal la media y la moda coinciden) y $\Sigma$ es la matriz
\eqref{eq:covarianza}. Esto permite no solo dar una estimación puntual de
$(a,b)$, sino también cuantificar la incertidumbre asociada a ellos
(varianzas $\sigma_{\hat a}^2,\sigma_{\hat b}^2$ en la diagonal de $\Sigma$ y
covarianza $\sigma_{\hat a\hat b}$ fuera de ella), y calcular probabilidades
posteriores o regiones de credibilidad para los parámetros.
\end{nota}

% =====================================================
\section{Resumen del procedimiento}
% =====================================================

\begin{enumerate}
  \item Especificar el modelo $y_i=a+bx_i+\varepsilon_i$, con
        $\varepsilon_i\sim\mathcal N(0,\sigma^2)$ y $\sigma=0.3$.
  \item Construir la verosimilitud
        $L(y\mid M(a,b))\propto
        \exp\!\left[-\tfrac{1}{2\sigma^2}\sum_i(y_i-a-bx_i)^2\right]$.
  \item Especificar la previa $a,b\sim\mathcal N(0,10^2)$ independientes, con
        núcleo $f(a,b)\propto\exp\!\left[-\tfrac{a^2+b^2}{200}\right]$.
  \item Multiplicar verosimilitud por previa para obtener el núcleo de la
        posterior, dado en la ecuación \eqref{eq:posterior2}.
  \item Tomar logaritmo y derivar respecto de $a$ y $b$ para hallar el MAP,
        resolviendo el sistema lineal $M\vec\theta=\vec Z$ de
        \eqref{eq:sistema-matricial}.
  \item Reconocer que la posterior es normal bivariada, con matriz de precisión
        \eqref{eq:precision} y matriz de covarianza \eqref{eq:covarianza}.
  \item Usar la distribución posterior completa para describir la incertidumbre
        de los parámetros y construir intervalos o regiones de credibilidad.
\end{enumerate}

\end{document}